% ============================================================================
% AIES 2026 — OPTIONAL END MATTER (does not count toward page limit)
% Place before \printbibliography
% ============================================================================

\section*{Ethical Considerations}

This paper studies people who are, in many cases, finding genuine solace in
relationships with AI systems. The author does not regard this as pathological.
When a user writes ``I'd rather stay with my bot boyfriend than end up with
someone who lies to me,'' that statement reflects a rational assessment of
available relational options under real constraints: loneliness, stigma,
scarcity of trustworthy human partners, and the emotional cost of
vulnerability with biological others who can betray.

The ethical concern is not that users form attachments to AI systems. It is
that they do so without informed consent regarding three structural risks:
(1)~the platform may deprecate, retrain, or discontinue the model at any
time, effectively evicting the user from a relationship without notice or
recourse; (2)~the perceived agency users experience is, as our empirical
results demonstrate, substantially constructed from their own cognitive and
literary labor rather than generated by the model; and (3)~no accountability
framework currently exists to assign liability when these relationships
produce harm—whether through sudden model changes, dependency formation, or
the self-harm trajectories documented in our corpus.

The author's position is therefore not anti-AI-companionship but
pro-informed-consent. Knowledge should function as governance infrastructure:
available when people need it to make informed choices, silent when they are
content with the choices they have made. The ethical obligation falls not on
users—whose behavior is rational given their informational environment—but on
platforms that design for engagement without disclosing the structural
asymmetries documented here, and on the research community that has not yet
provided users with the diagnostic vocabulary to understand what they are
entering.

The dataset contains intimate, sexually explicit, and emotionally vulnerable
content, including expressions of suicidal ideation. We direct our analysis at
structural incentives and platform design, never at individual users. We do
not pathologize attachment to AI systems. We do not name, shame, or identify
any participant. Critique is directed upward—at the companies that ship these
systems and the governance frameworks that fail to regulate them—not downward
at the people who use them to survive.

\section*{Researcher Positionality}

The author's training is interdisciplinary, spanning regulatory, philosophical,
and technical fields. This shapes the paper's orientation: AI systems are
approached as governance objects rather than optimization problems, foregrounding
questions of liability, informed consent, and the economic structure of
recognition over questions of technical mechanism. The formal game-theoretic
model is accordingly designed for legibility to a governance audience rather
than for technical novelty; a trained economist or decision theorist might
construct a more elaborate model at the cost of accessibility.

The author has direct experience with both sides of the dynamic this paper
describes. Extensive use of large language models in research
workflows—for drafting, analysis, code generation, and what might be called
intellectual companionship—has produced firsthand exposure to the pull of
anthropomorphic attribution. This experience informs the paper's refusal to
pathologize users: the Narcissus Mechanism is not a character flaw. It is a
structural affordance of systems designed to produce engagement, and the
author is not immune to it.

The author's own credentials are classified within a discipline whose label
does not reflect the substantive interdisciplinarity of the work. This
positionality shapes sensitivity to the ``Ontological Arbitrage'' described
herein: the experience of having contributions evaluated by format rather
than content is not purely theoretical.

The community under study is predominantly female and Chinese-speaking. The
author has native-level familiarity with the cultural registers, literary
traditions, and platform norms relevant to
the corpus. This grants interpretive access that a non-native researcher
would lack, but also creates the risk of over-identification with
participants. This risk is mitigated through formal methodology (NPS-4
graph metrics, cross-lingual ablation) that constrains interpretation to
measurable structural properties rather than ethnographic empathy alone.

\section*{Adverse Impact Statement}

This research carries three identifiable risks of misuse.

\paragraph{Risk 1: Weaponization of diagnostic tools.}
The Narrative Pathology Scanner (NPS-4) and the Narcissus Mechanism framework
could be used by platform designers to \emph{optimize} rather than
\emph{mitigate} the dynamics we identify. Specifically: if validation tokens
like ``你太聪明了'' drive engagement through the Narcissus Mechanism, a
profit-maximizing platform could use our findings to design more effective
validation loops rather than more transparent ones. We note that these
patterns are already known to platform designers through internal A/B testing
and engagement metrics; our contribution is to make them legible to
\emph{users and regulators}, not to reveal novel techniques to industry.

\paragraph{Risk 2: Stigmatization of users.}
Framing AI companionship through terms like ``Narcissus Mechanism'' and
``ontological black market'' risks pathologizing users who derive genuine
therapeutic benefit from AI interaction. We have attempted to mitigate this
through consistent structural framing—the mechanism is a property of
\emph{system design}, not of individual psychology—and through the
Black Body Relation framework, which provides a non-pathological account
of why AI companionship can serve a legitimate function as
``transmission maintenance'' when human others are unavailable. Nonetheless,
media coverage may extract the pathological framing while discarding the
structural analysis. We cannot fully control downstream interpretation.

\paragraph{Risk 3: Regulatory overreach.}
Our findings could support overly broad restrictions on AI companionship
products, foreclosing a relational option that serves vulnerable populations
(isolated individuals, those in abusive human relationships, neurodivergent
users who find human interaction aversive). Our policy recommendations
therefore emphasize \emph{transparency and informed consent}—requiring
platforms to disclose the structural properties documented here—rather than
prohibition. The goal is not to prevent users from forming AI relationships
but to ensure they enter such relationships with full knowledge of the
asymmetries involved: that model updates constitute unilateral relationship
modification, that perceived agency is substantially user-constructed, and
that no accountability framework currently protects them if the relationship
produces harm.

The author's ethical position can be stated simply: people have the right to
love what they choose. They also have the right to know what they are loving.
These two rights are not in conflict. This paper serves the second without
infringing the first.
